\documentclass[11pt]{article}
\usepackage[a4paper,margin=27mm]{geometry}
\usepackage{amsmath,amssymb,mathtools}
\usepackage{fontspec,unicode-math}
\setmainfont{Latin Modern Roman}
\setmathfont{Latin Modern Math}
\usepackage{xurl}
\usepackage[hidelinks]{hyperref}
\newcommand{\sref}[2]{\hyperlink{source:#1:#2}{[S#2]}}
\newcommand{\cataloguescope}[1]{\paragraph{Mathematical question.}#1}
\hypersetup{pdftitle={Problem statement}}
\begin{document}
\section{Do any (2, 5)-perfect numbers exist}\label{statement}

\subsection{Definitions and mathematical statement}

\paragraph{Shared notation.}$\mathbb N=\{1,2,\ldots\}$, and $\mathbb Z,\mathbb Q,\mathbb R,\mathbb C$ denote the integers, rationals, reals and complex numbers. Any other conventions are specified locally.


For $n\in\mathbb N$, write $\sigma(n)=\sum_{d\mid n,\ d>0}d$. The notation $(m,k)$-perfect means that $m$ successive applications of $\sigma$ multiply the original integer by $k$: $\sigma^{\circ m}(n)=kn$. In particular, $(2,5)$-perfect means $\sigma(\sigma(n))=5n$.
\paragraph{Statement recorded in the supplied source.}
Does there exist $n\in\mathbb N$ such that
\[\sigma(\sigma(n))=5n?\] \sref{statement}{2}
\subsection{Short English statement}
Does any positive integer have a sum of divisors whose own sum of divisors is five times the original integer?
\subsection{Sources}
\begin{description}
\item[\hypertarget{source:statement:1}{S1}] ULAM.AI and the UnsolvedMath Contributors, UnsolvedMath: A Curated Collection of Open Mathematics Problems, catalogue entry. CC BY 4.0. \url{https://huggingface.co/datasets/ulamai/UnsolvedMath}.
\item[\hypertarget{source:statement:2}{S2}] Wikipedia contributors, \emph{List of unsolved problems in mathematics}, revision 1366636767 (2026) Number theory, Do any (2, 5)-perfect numbers exist. \url{https://en.wikipedia.org/w/index.php?title=List_of_unsolved_problems_in_mathematics&oldid=1366636767#Number_theory}.
\item[\hypertarget{source:statement:3}{S3}] Wikipedia contributors, \emph{Superperfect number}, revision 1365374288 (2026) Generalizations: (m,k)-perfect numbers. \url{https://en.wikipedia.org/w/index.php?title=Superperfect+number&oldid=1365374288}.
\end{description}
\label{end:statement}
\end{document}
